\subsection{Constraints on Large Extra Dimensions from the MINOS Experiment --- arXiv:1608.06964}

\textbf{Authors:} MINOS Collaboration, P. Adamson et al.

\paragraph{主な主張}
MINOSの実データにLEDの有意な兆候は見られず、最軽ニュートリノ質量がゼロの極限で$R<0.45\,\mu\mathrm{m}$という90\%信頼水準の上限を得る\cite{MINOS:2016LEDConstraints}。

\paragraph{新規性}
近検出器・遠検出器の荷電カレントおよび中性カレント事象のエネルギー分布を用い、近検出器での振動も含めた実験コラボレーション自身のLED探索を行った。

\paragraph{位置付け}
将来感度の予測ではなく実測に基づく長基線制約であり、sterileへの消失を検出器応答・系統誤差込みで検証する解析例である。
